\chapter{CONCLUSIONS}\label{ch:conclusion}
The exercise demonstrated that a simple visual system significantly enhances workflow transparency and team coordination. Parallel execution of independent tasks increased throughput, but ambiguities in task description led to rework and delays, underscoring the need for precise initial requirements. The built-in "Test" stage proved crucial for quality assurance, catching inconsistencies. Key learnings include the importance of clear task definition, the efficiency gained from role specialization, and the value of a visible workflow for identifying and resolving bottlenecks in real-time. The Kanban method is highly effective for managing collaborative, step-by-step processes.